$-- A single line per entry: the date at the left, and beside it the entry
$-- itself, with `what` picked out in bold as upstream's tables do.
$--
$-- The first loop measures the dates so the column fits them exactly, the way
$-- a tabular would; the second sets the entries. Neither the document nor the
$-- author has a width to choose.
\begin{cvlisting}
$for(entries)$\cvmeasureentry{$entries.when$}$endfor$
$for(entries)$
\cventryline{$entries.when$}{$if(entries.what)$\textbf{$entries.what$}$endif$$if(entries.with)$$if(entries.what)$, $endif$$entries.with$$endif$$if(entries.where)$, \emph{$entries.where$}$endif$}
$endfor$
\end{cvlisting}
